%%%%%%%%%%%%%%%%%%%%%%%%%%%%%%%%%%%%%%%%%%%%%%%%%%%%%%%%%%%%%%%%%%%%%%
% Lớp 11 | Hàm số lượng giác và phương trình lượng giác | Phương trình lượng giác cơ bản
% Dạng 1D1?5-4: Phương trình cơ bản dùng Độ
% Số câu: 24  (sinh tự động bởi tach_cau.py; từ nay thầy sửa trực tiếp file này)
%%%%%%%%%%%%%%%%%%%%%%%%%%%%%%%%%%%%%%%%%%%%%%%%%%%%%%%%%%%%%%%%%%%%%%

%%%%%%%%%%%%%%%% TRẮC NGHIỆM 4 PHƯƠNG ÁN %%%%%%%%%%%%%%%%

% nguồn: Ok_Toan11_TracNghiem_Lan2
\begin{ex}%[1D1H5-4]%[ma:2a70959f]
	Tổng các nghiệm của phương trình $2\sin (x  + 40^\circ) = \sqrt{3}$ trên khoảng $(-180^\circ; 180^\circ)$ là
	 \choice
		{$20^\circ$} 
		{\True $100^\circ$} 
		{$80^\circ$} 
		{$120^\circ$} 
	 \loigiai{
		Ta  có
		\begin{align*}
			2\sin (x +  40^\circ) = \sqrt{3} &\Leftrightarrow \sin (x + 40^\circ) = \dfrac{\sqrt{3}}{2} \Leftrightarrow \sin (x + 40^\circ) = \sin60^\circ\\
			& \Leftrightarrow \hoac{ 
			 & x + 40^\circ = 60^\circ + k360^\circ \\
			& x + 40^\circ = 120^\circ  + k360^\circ} \Leftrightarrow \hoac{
			&  x = 20^\circ + k360^\circ \\
			& x = 80^\circ + k360^\circ}, k \in \mathbb{Z}. 
		\end{align*}
		Theo đề  bài
		\begin{itemize}
			\item $-180^\circ < 20^\circ + k360^\circ < 180^\circ  \Leftrightarrow -\dfrac{5}{9} < k < \dfrac{4}{9} \Rightarrow k = 0 \Rightarrow x = 20^\circ$.
			\item $-180^\circ < 80^\circ + k360^\circ <  180^\circ \Leftrightarrow -\dfrac{13}{18} < k < \dfrac{5}{18} \Rightarrow k = 0 \Rightarrow x = 80^\circ$.
		\end{itemize}
		 Vậy tổng các nghiệm của phương trình là $20^\circ + 80^\circ = 100^\circ$.
	 }
\end{ex}

% nguồn: Ok_Toan11_TracNghiem_Lan2
\begin{ex}%[1D1H5-4]%[ma:1ed917c6]
	Phương trình $\sin x = \sin  80^\circ$ có tập nghiệm là
	 \choice
		{\True $x=80^\circ +  k360^\circ$; $x=100^\circ + k360^\circ$, $k \in \mathbb{Z}$}
		{$x=80^\circ + k360^\circ$; $x=-80^\circ  + k360^\circ$, $k \in \mathbb{Z}$}
		 {$x=40^\circ + k360^\circ$; $x=140^\circ + k360^\circ$, $k \in \mathbb{Z}$}
		{$x=80^\circ + k180^\circ$;  $x=100^\circ + k180^\circ$, $k \in \mathbb{Z}$}
	\loigiai{ 
		Ta có \begin{align*}
		\sin x = \sin 80^\circ 
		 & \Leftrightarrow \hoac{
		& x =  80^\circ + k360^\circ \\
		& x = 180^\circ - 80^\circ + k360^\circ} 
		\Leftrightarrow \hoac{ 
		& x = 80^\circ  + k360^\circ \\
		& x = 100^\circ  + k360^\circ} \, ,k \in \mathbb{Z}.
	\end{align*}
	Vậy tập nghiệm của phương trình là $x=80^\circ + k360^\circ$; $x=100^\circ +  k360^\circ$, $k \in \mathbb{Z}$.
} 
\end{ex}

% nguồn: Ok_Toan11_TracNghiem_Lan2
\begin{ex}%[1D1H5-4]%[ma:56ac069d]
	Nghiệm của phương trình $\sin x =  \sin 30^\circ$ là
	 \choice
		{\True $x= 30^\circ +  k360^\circ$; $x = 150^\circ + k360^\circ$, $k \in \mathbb{Z}$}
		 {$x = \pm 30^\circ + k2\pi$, $k \in \mathbb{Z}$}
		{$x=\pm 30^\circ + k360^\circ$, $  k \in \mathbb{Z}$}
		{$x=30^\circ + 360^\circ$; $x=150^\circ + 360^\circ$, $k \in \mathbb{Z}$} 
	\loigiai{ 
		Ta  có
		\begin{align*}
			\sin  x = \sin 30^\circ \Leftrightarrow \hoac{
			 & x = 30^\circ + k360^\circ \\
			& x = 180^\circ -  30^\circ + k360^\circ} \Leftrightarrow \hoac{
			& x =  30^\circ + k360^\circ \\
			 & x = 150^\circ + k360^\circ}, k \in \mathbb{Z}.
		\end{align*}
		Vậy nghiệm của phương trình là $x= 30^\circ + k360^\circ$; $x=150^\circ +  k360^\circ$, $ k \in \mathbb{Z}$.
	 }
\end{ex}

% nguồn: Ok_Toan11_TracNghiem_Lan2
\begin{ex}%[1D1H5-4]%[ma:99ec5964]
	Tập nghiệm của phương trình  $\sin x=\sin(x-60^\circ)$ là
	\choice 
		{$\left\{\dfrac{\pi}{3}+k\pi\,\mid  dle|\ k\in\mathbb{Z}\right\}$}
		 {$\left\{\dfrac{2\pi}{3}+k\pi\,\mid dle|\ k\in\mathbb{Z}\right\}$}
		{\True $\left\{120^\circ+k180^\circ\,\mid dle|\  k\in\mathbb{Z}\right\}$}
		{$\left\{60^\circ+k180^\circ\,\mid  dle|\ k\in\mathbb{Z}\right\}$}
	\loigiai{ 
		Ta có 
		\allowdisplaybreaks 
		\begin{eqnarray*}
			 &&\sin x=\sin(x-60^\circ)\\
			&\Leftrightarrow&\hoac{ 
			& x=x-60^\circ+k360^\circ\\ 
			& x=180^\circ-(x-60^\circ)+k360^\circ 
			 }\\
			&\Leftrightarrow& x=120^\circ+k180^\circ~  (k\in\mathbb{Z}).
		\end{eqnarray*}
	} 
\end{ex}

% nguồn: Ok_Toan11_TracNghiem_Lan2
\begin{ex}%[1D1H5-4]%[ma:46e103c6]
	Nghiệm của phương trình $\cos x=\cos  60^{\circ}$ là
	\choice 
		{$x=120^{\circ}+k180^{\circ}, k\in\mathbb{Z}$} 
		{$\hoac{&x=60^{\circ}+k360^{\circ}\\ 
		&x=120^{\circ}+k360^{\circ}}  (k\in\mathbb{Z})$}
		{\True $x=\pm 60^{\circ}+k360^{\circ}, k\in\mathbb{Z}$} 
		{$x=60^{\circ}+k180^{\circ}, k\in\mathbb{Z}$} 
	\loigiai{ 
		Ta có $\cos  x=\cos 60^{\circ}\Leftrightarrow x=\pm 60^{\circ}+k360^{\circ}, k\in\mathbb{Z}$.\\
		Vậy  nghiệm của phương trình đã cho là $x=\pm 60^{\circ}+k360^{\circ}, k\in\mathbb{Z}$.
	 }
\end{ex}

% nguồn: Ok_Toan11_TracNghiem_Lan2
\begin{ex}%[1D1H5-4]%[ma:d811c6bb]
	Tổng các nghiệm phương trình $\tan \left( 2x-15^\circ \right)=1$ trên khoảng $\left( -90^\circ;90^\circ  \right)$ bằng
	 \choice
		 {$30^\circ$}
		 {$-60^\circ$}
		 {$0^\circ$}
		{\True  $-30^\circ$}
	\loigiai{Ta có phương trình $\tan \left( 2x-15^\circ \right)=1\Leftrightarrow 2x-15^\circ=45^\circ+k180^\circ\Leftrightarrow x=30^\circ+k90^\circ$ 
		với  $k\in \mathbb{Z}$. \\
		Xét $-90^\circ<30^\circ+k90^\circ<90^\circ\Leftrightarrow -\dfrac{4}{3}<k<\dfrac{2}{3}\Rightarrow \hoac{&k=-1 \\&k=0 }\Rightarrow \hoac{&x_1=-60^\circ  \\ &x_2=30^\circ }$. \\
	 Vậy $x_1+x_2=-30^\circ$.}
\end{ex}

% nguồn: Ok_Toan11_TracNghiem_Lan2
\begin{ex}%[1D1H5-4]%[ma:47b85926]
	Nghiệm  của phương trình $2\cos \left(x-15^{\circ}\right)-1=0$ là
	\choice 
		{$\hoac{& x=75^{\circ}+k 360^{\circ} \\ & x=135^{\circ}+k 360^{\circ}}$, $k \in  \mathbb{Z}$}
		{$\hoac{& x=60^{\circ}+k 360^{\circ} \\ & x=-60^{\circ}+k 360^{\circ}}$, $k  \in \mathbb{Z}$}
		{$\hoac{& x=45^{\circ}+k 360^{\circ} \\ &  x=-45^{\circ}+k 360^{\circ}}$, $k \in \mathbb{Z}$}
		{\True $\hoac{& x=75^{\circ}+k 360^{\circ} \\ & x=-45^{\circ}+k 360^{\circ}}$,  $k \in \mathbb{Z}$}
	\loigiai{ 
		Ta  có
		\allowdisplaybreaks 
		\begin{eqnarray*}
			 &&2 \cos \left(x-15^{\circ}\right)-1=0 \\
			 &\Leftrightarrow& \cos \left(x-15^{\circ}\right)=\dfrac{1}{2}\\ &\Leftrightarrow& \cos \left(x-15^{\circ}\right)=\cos 60^{\circ}\\
			 &\Leftrightarrow& \hoac{x-15^{\circ}&=60^{\circ}+k360^{\circ} \\ x-15^{\circ}&=-60^{\circ}+k 360^{\circ}}\\
			&\Leftrightarrow& \hoac{& x=75^{\circ}+k 360^{\circ} \\ & x=-45^{\circ}+k 360^{\circ}},\,  k \in \mathbb{Z}.
		\end{eqnarray*}
	 }
\end{ex}

% nguồn: Ok_Toan11_TracNghiem_Lan2
\begin{ex}%[1D1H5-4]%[ma:9e5b5242]
	Tìm nghiệm của phương trình $\cot  x=\sqrt{3}$.
	\choice 
		 {$x=60^{\circ}+k180^{\circ}$}
		{$x=60^{\circ}+k360^{\circ}$} 
		{\True $x=30^{\circ}+k180^{\circ}$} 
		{$x=-60^{\circ}+k180^{\circ}$} 
	 \loigiai{
		 Ta có $\cot x=\sqrt{3}\Leftrightarrow\cot x=\cot 30^{\circ}\Leftrightarrow x=30^{\circ}+k180^{\circ}, k\in\mathbb{Z}$.\\
		Vậy nghiệm của phương trình đã cho là $x=30^{\circ}+k180^{\circ},  k\in\mathbb{Z}$.
	 }
\end{ex}

% nguồn: Ok_TracNghiemLop11_Lan1
\begin{ex}%[1D1H5-4]%[ma:7b1264b0]
	Phương trình $\cot \left(4x-20^\circ  \right)=\dfrac{1}{\sqrt{3}}$ có họ nghiệm là
	\choice 
		{$x=30^\circ+k\,45^\circ,  k \in \mathbb{Z}$}
		{$x=20^\circ+k\,90^\circ,  k \in \mathbb{Z}$}
		 {$x=35^\circ+k\,90^\circ, k \in \mathbb{Z}$}
		{\True $x=20^\circ+k\,45^\circ, k \in \mathbb{Z}$} 
	\loigiai{ 
		 Ta có
		\begin{eqnarray*}
			& & \cot \left(4x-20^\circ \right)=\dfrac{1}{\sqrt{3}}\\ 
			&\Leftrightarrow&  4x-20^\circ=60^\circ+k\,180^\circ\\
			&\Leftrightarrow& 4x=80^\circ+k\,180^\circ\\ 
			&\Leftrightarrow& x=20^\circ+k\,45^\circ, k \in \mathbb{Z}. 
		\end{eqnarray*}
	} 
\end{ex}

% nguồn: Ok_TracNghiemLop11_Lan1
\begin{ex}%[1D1H5-4]%[ma:17e2206b]
	Giải phương trình $\cot 2x=\cot  20^\circ $.
	\choice 
		{\True $x=10^\circ +k90^\circ ,k\in  \mathbb{Z}$}
		{$x=10^\circ +k180^\circ ,k\in \mathbb{Z}$} 
		{$x=20^\circ +k90^\circ ,k\in \mathbb{Z}$} 
		{$x=20^\circ +k180^\circ  ,k\in \mathbb{Z}$}
	\loigiai{ 
		Ta có $\cot 2x=\cot 20{}^\circ \Leftrightarrow 2x=20^\circ  +k180^\circ \Leftrightarrow x=10{}^\circ +k90^\circ ,k\in \mathbb{Z}$.
	} 
\end{ex}

% nguồn: Ok_TracNghiemLop11_Lan1
\begin{ex}%[1D1H5-4]%[ma:2e22cfae]
	Số nghiệm của phương trình  $\sin 3x = 0$ với $0 < x < 180^\circ$ là
	\choice 
		{$1$} 
		{$3$} 
		 {\True $2$}
		 {$4$}
	\loigiai{ 
	Ta có: $\sin 3x = 0 \Leftrightarrow x = \dfrac{k  180^\circ}{3} \quad(k \in \mathbb{Z})$.\\
	Do $0 <  x < 180^\circ$ nên ta có $0 < \dfrac{k 180^\circ}{3} < 180^\circ \Leftrightarrow 0 < k < 3$.\\
	Do $k \in \mathbb{Z}$ nên $k  \in \lbrace 1;2 \rbrace$.\\
	Với $k=1$ ta được nghiệm $x  = 60^\circ $.\\
	Với  $k=2$ ta được nghiệm $x = 120^\circ$.\\
	Vậy  có $2$ nghiệm thỏa mãn.
	 }
\end{ex}

% nguồn: Ok_TracNghiemLop11_Lan1
\begin{ex}%[1D1H5-4]%[ma:85a8dbef]
	Tìm nghiệm của phương trình  $\cos (x+30^\circ)=-\dfrac{\sqrt{3}}{2}$.
	 \choice
		{$x=120^\circ (1+3k)$ hoặc $x=180^\circ (k-1)$, $k\in \mathbb Z$} 
		{\True $x=120^\circ  +k360^\circ$ hoặc $x=-180^\circ +k360^\circ$, $k\in \mathbb Z$}
		{$x=120^\circ +k180^\circ$ hoặc $x=-180^\circ +k180^\circ$, $k\in  \mathbb Z $}
		 {$x=120^\circ +k2\pi$ hoặc $x=-180^\circ +k2\pi$, $k\in \mathbb Z $}
	 \loigiai{
		$\cos (x+30^\circ)=-\dfrac{\sqrt{3}}{2} \Leftrightarrow \hoac{&x+30^\circ=150^\circ+k360^\circ\\&x+30^\circ=-150^\circ+k360^\circ} \Leftrightarrow \hoac{&x=120^\circ+k360^\circ\\&x=-180^\circ+k360^\circ} (k\in \mathbb  Z)$.
	} 
\end{ex}

% nguồn: Ok_TracNghiemLop11_Lan1
\begin{ex}%[1D1H5-4]%[ma:008fd1fc]
	Giải phương trình $\tan (2x)=\tan 80^\circ $. Kết quả thu được là 
	\choice 
		 {$x=80^\circ +k180^\circ $}
		{\True  $x=40^\circ +k90^\circ $}
		 {$x=40^\circ +k45^\circ $}
		{$x=40^\circ +k180^\circ $} 
	\loigiai{ 
		Ta có $\tan \left( 2x  \right)=\tan 80^\circ $$\Leftrightarrow 2x=80^\circ +k180^\circ $$\Leftrightarrow x=40^\circ +k90^\circ ,k\in \mathbb{Z}$.
	 }
\end{ex}

% nguồn: Ok_Toan11_TracNghiem_Lan2
\begin{ex}%[1D1V5-4]%[ma:10255d24]
	Tổng các nghiệm của phương trình $ 2\sin\left(x+40^\circ\right)=\sqrt{3}$ trên  khoảng $\left(-180^\circ\,;\,180^\circ\right)$ là
	\choice 
		{$20^\circ$} 
		{\True  $100^\circ$}
		 {$80^\circ$}
		{$120^\circ$} 
	 \loigiai{
		Ta có 
		$$2\sin\left(x+40^\circ\right)=\sqrt{3}\Leftrightarrow\sin\left(x+40^\circ\right)=\dfrac{\sqrt{3}}{2}\Leftrightarrow\hoac{& x+40^\circ=60^\circ+k360^\circ\\ & x+40^\circ=120^\circ+k360^\circ}\,(k\in\mathbb{Z})\Leftrightarrow\hoac{&x=20^\circ+k360^\circ\\&  x=80^\circ+k360^\circ}\,(k\in\mathbb{Z}).$$
		Xét  $x\in \left(-180^\circ\,;\,180^\circ\right)$, ta có
		\begin{itemize}
			\item $-180^\circ <20^\circ+k360^\circ <180^\circ\Leftrightarrow-\dfrac{5}{9}<k<\dfrac{4}{9}\Rightarrow k=0\Rightarrow x=20^\circ  $.
			\item $-180^\circ <80^\circ+k360^\circ <180^\circ\Leftrightarrow-\dfrac{13}{18}<k<\dfrac{5}{18}\Rightarrow k=0\Rightarrow x=80^\circ$. 
		\end{itemize}
		Vậy tổng  các nghiệm của phương trình là $20^\circ+80^\circ=100^\circ$.
	 }
\end{ex}


%%%%%%%%%%%%%%%% ĐÚNG/SAI %%%%%%%%%%%%%%%%

% nguồn: Ok_dungsai_lop11_Lan1
\begin{ex}%[1D1N5-4]%[ma:9f360346]
 Xét tính đúng sai của các  mệnh đề sau.
  \choiceTF
  {$\sin  x=\sin 15^{\circ}\Leftrightarrow\Leftrightarrow\hoac{&x=15^{\circ}+k360^{\circ}\\&x=165^{\circ}+k360^{\circ}}; k\in\mathbb{Z}$}
  {$\sin x=\sin 15^{\circ}\Leftrightarrow x=\pm  15^{\circ}+k360^{\circ}; k\in\mathbb{Z}$}
  {\True $\sin x=\sin 15^{\circ}\Leftrightarrow\sin  x=\cos75^\circ$}
   {\True $\sin x=\sin 15^{\circ}\Leftrightarrow\hoac{&x=15^{\circ}+k360^{\circ}\\&x=165^{\circ}+k360^{\circ}}; k\in\mathbb{Z}$}
 \loigiai{\begin{itemize}
 \item Ta có  $\sin x=\sin 15^{\circ}\Leftrightarrow\sin x=\cos75^\circ$ vì $\cos75^\circ=\sin15^\circ$.
 \item $\sin x=\sin 15^{\circ}\Leftrightarrow\hoac{&x=15^{\circ}+k360^{\circ}\\&x=180^{\circ}-15^{\circ}+k360^{\circ}}\Leftrightarrow\hoac{&x=15^{\circ}+k360^{\circ}\\&x=165^{\circ}+k360^{\circ}}; k\in\mathbb{Z}$. 
\end{itemize}}
\end{ex}

% nguồn: Ok_dungsai_lop11_Lan1
\begin{ex}%[1D1N5-4]%[ma:f114553d]
 Xét tính đúng sai của các mệnh đề  sau.
  \choiceTF
  {$\sin  x=\sin 15^{\circ}\Leftrightarrow x=\pm 15^{\circ}+k360^{\circ}; k\in\mathbb{Z}$}
  {\True  $\sin x=\sin 15^{\circ}\Leftrightarrow\sin x=\cos75^\circ$}
  {$\sin x=\sin 15^{\circ}\Leftrightarrow\Leftrightarrow\hoac{&x=15^{\circ}+k360^{\circ}\\&x=165^{\circ}+k360^{\circ}};  k\in\mathbb{Z}$}
  {\True  $\sin x=\sin 15^{\circ}\Leftrightarrow\hoac{&x=15^{\circ}+k360^{\circ}\\&x=165^{\circ}+k360^{\circ}}; k\in\mathbb{Z}$}
 \loigiai{\begin{itemize}
 \item Ta  có $\sin x=\sin 15^{\circ}\Leftrightarrow\sin x=\cos75^\circ$ vì $\cos75^\circ=\sin15^\circ$.
  \item $\sin x=\sin 15^{\circ}\Leftrightarrow\hoac{&x=15^{\circ}+k360^{\circ}\\&x=180^{\circ}-15^{\circ}+k360^{\circ}}\Leftrightarrow\hoac{&x=15^{\circ}+k360^{\circ}\\&x=165^{\circ}+k360^{\circ}}; k\in\mathbb{Z}$.
\end{itemize}}
\end{ex}

% nguồn: Ok_dungsai_lop11_Lan1
\begin{ex}%[1D1N5-4]%[ma:89d86688]
 Tìm công  thức nghiệm (theo đơn vị độ) của phương trình $\cot x=-1$, với ký hiệu $k\in\mathbb{Z}$.
 \choiceTF 
  {$x=45^{\circ}+k180^{\circ}$} 
  {\True  $x=-45^{\circ}+k180^{\circ}$}
  {$x=45^{\circ}+k360^{\circ}$} 
  {\True  $x=135^{\circ}+k180^{\circ}$}
 \loigiai{ 
   Ta có
  $\cot x=1\Leftrightarrow\cot x=\cot\left( -45^{\circ}\right)\Leftrightarrow  x=-45^{\circ}+k180^{\circ}=135^{\circ}+k180^{\circ} (k\in\mathbb{Z}).$
 } 
\end{ex}

% nguồn: Ok_dungsai_lop11_Lan1
\begin{ex}%[1D1H5-4]%[ma:79cb88c6]
 Cho phương trình lượng giác $\tan (2 x-15^{\circ})=1\qquad (*)$. Khi  đó
 \choiceTF 
   {\True Phương trình $(*)$ có nghiệm $x=30^{\circ}+k 90^{\circ},\;(k \in \mathbb{Z})$}
  {Phương trình có nghiệm âm lớn nhất bằng  $-30^{\circ}$}
  {Tổng các  nghiệm của phương trình trong khoảng $(-180^{\circ}; 90^{\circ})$ bằng $180^{\circ}$}
  {Trong khoảng $(-180^{\circ}; 90^{\circ})$ phương trình có nghiệm lớn nhất  bằng $60^{\circ}$}
 \loigiai{ 
  \begin{itemchoice}
   \itemch \textbf{Đúng}. Ta có \begin{eqnarray*}
    &&\tan (2x-15^\circ)=1\\
   &\Leftrightarrow&2x-15^\circ=45^\circ+k\cdot 180^\circ\\ 
   &\Leftrightarrow&2x=60^\circ+k\cdot  180^\circ\\
    &\Leftrightarrow&x=30^\circ+k\cdot 90^\circ,k\in \mathbb{Z}.
  \end{eqnarray*}
  Vậy phương trình $(*)$ có nghiệm $x=30^\circ+k\cdot 90^\circ,k\in \mathbb{Z}$. 
  \itemch \textbf{Sai}. Xét $x=30^\circ+k\cdot 90^\circ<0\Leftrightarrow k<-\dfrac{1}{3}$. Chọn  $k=-1$.\\
  Khi  đó phương trình có nghiệm âm lớn nhất bằng $x=30^\circ-90^\circ=-60^\circ$.
  \itemch \textbf{Sai}. Xét $k\in \mathbb{Z}$ và 
  \begin{eqnarray*}
    &&-180^\circ<30^\circ+k\cdot 90^\circ<90^\circ\\
   &\Leftrightarrow&-210^\circ<k\cdot  90^\circ<60^\circ\\
   &\Leftrightarrow&-\dfrac{7}{3}<k<\dfrac{2}{3}\\ 
    &\Leftrightarrow&k\in \{-2;-1;0\}.
  \end{eqnarray*}
  Suy ra các nghiệm  là $\{-150^\circ; -60^\circ; 30^\circ\}$.\\
  Vậy tổng các nghiệm là  $(-150^\circ) + (-60^\circ) + 30^\circ=-180^\circ$.
  \itemch  \textbf{Sai}. Trong khoảng $(-180^{\circ}; 90^{\circ})$ phương trình có nghiệm lớn nhất bằng $30^\circ$.
 \end{itemchoice}
} 
\end{ex}

% nguồn: Ok_dungsai_lop11_Lan2
\begin{ex}%[1D1H5-4]%[ma:dac917d2]
	Cho phương trình lượng giác  $\tan\left(2x-15^\circ\right)=1$. Khi đó
	\choiceTF 
		{\True Phương trình đã  cho có nghiệm $x=30^\circ+k\cdot 90^\circ\,\left(k\in\mathbb{Z}\right)$}
		{Trong khoảng $\left(-180^\circ;90^\circ\right)$ thì phương trình có  nghiệm âm lớn nhất bằng $-30^\circ$}
		 {Tổng các nghiệm của phương trình trong khoảng $\left(-180^\circ;90^\circ\right)$ bằng $180^\circ$}
		{Trong  khoảng $\left(-180^\circ;90^\circ\right)$ phương trình có nghiệm lớn nhất bằng $60^\circ$}
	 \loigiai{
		\begin{itemchoice}
			\itemch Ta có $\tan\left(2x-15^\circ\right)=1\Leftrightarrow  x=30^\circ+k\cdot 90^\circ\,(k\in\mathbb{Z})$.
			\itemch  Ta có $-180^\circ <x<90^\circ\Rightarrow-180^\circ <30^\circ+k90^\circ <90^\circ\, (k\in\mathbb{Z})\Rightarrow k=\{-2;-1;0\}\Rightarrow\hoac{&x=-150^\circ\\&x=-60^\circ\\&x=30^\circ.}$\\
			Vậy trong khoảng  $\left(-180^\circ;90^\circ\right)$ thì phương trình có nghiệm âm lớn nhất bằng $-60^\circ$.
			\itemch Tổng các nghiệm của phương trình trong khoảng $\left(-180^\circ;90^\circ\right)$ bằng  $$-150^\circ+\left(-60\right)^\circ+30^\circ=-180^\circ.$$
			\itemch Trong khoảng $\left(-180^\circ;90^\circ\right)$ phương trình có nghiệm lớn nhất bằng  $30^\circ $.
		\end{itemchoice}
	} 
\end{ex}

% nguồn: Ok_dungsai_lop11_Lan2
\begin{ex}%[1D1V5-4]%[ma:ee3d75ca]
	 Cho phương trình lượng giác $\sqrt{2}-2\sin\left(45^{\circ}-2x\right)=0 \quad (*)$.
	 \choiceTF
		{\True Phương trình $(*)$ tương đương  với $\sin\left(45^{\circ}-2x\right)=\sin 45^{\circ}$}
		 {\True Đồ thị hàm số $y=\sqrt{2}-2\sin\left(45^{\circ}-2x\right)$ cắt trục hoành tại gốc tọa độ}
		{\True Phương trình $(*)$ có nghiệm là $x=-k180^{\circ}$; $x=-45^{\circ}-k180^{\circ}, (k\in\mathbb{Z})$} 
		{Trên khoảng $\left(-\dfrac{\pi}{2};\dfrac{\pi}{2}\right)$  phương trình $(*)$ có một nghiệm}
	 \loigiai{
		Ta có 
		\begin{eqnarray*}
			 &&\sqrt{2}-2\sin\left(45^{\circ}-2x\right)=0
			 \Leftrightarrow\sin\left(45^{\circ}-2x\right)=\dfrac{\sqrt{2}}{2}
			\Leftrightarrow\sin\left(45^{\circ}-2x\right)=\sin  45^{\circ} \\
			&\Leftrightarrow&\hoac{&45^{\circ}-2x=45^{\circ}+k360^{\circ}\\&45^{\circ}-2x=180^{\circ}-45^{\circ}+k360^{\circ}},(k\in\mathbb{Z}) 
			 \Leftrightarrow\hoac{&x=-k180^{\circ}\\&x=-45^{\circ}-k180^{\circ}},(k\in\mathbb{Z}).
		\end{eqnarray*}
		\begin{itemchoice}
			\itemch  Phương trình tương đương với $\sin\left(45^{\circ}-2x\right)=\sin 45^{\circ}$.
			\itemch Ta có $0=\sqrt{2}-2\sin\left(45^{\circ}-2\cdot0\right)$ nên đồ thị  hàm số đi qua gốc tọa độ.
			\itemch Vậy phương trình có nghiệm là  $x=-k180^{\circ};x=-45^{\circ}-k180^{\circ},(k\in\mathbb{Z})$.
			\itemch Trên  khoảng $\left(-\dfrac{\pi}{2};\dfrac{\pi}{2}\right)$ phương trình đã cho có hai nghiệm.
		\end{itemchoice}
	} 
\end{ex}


%%%%%%%%%%%%%%%% TRẢ LỜI NGẮN %%%%%%%%%%%%%%%%

% nguồn: Ok_TraLoiNgan_Lop11_Lan1
\begin{ex}%[1D1H5-4]%[ma:0ad5e554]
	Phương trình $\tan \left(x+30^\circ\right)+1=0$ có bao  nhiêu nghiệm thỏa mãn $-90^\circ<x<360^\circ$?
	 \shortans{3}
	\loigiai{ 
		Ta có: $\tan \left(x+30^\circ\right)+1=0\Leftrightarrow \tan  \left(x+30^\circ\right)=-1=\tan\left(-45^\circ\right)\Leftrightarrow x=-75^\circ+k180^\circ$.\\
		Do $-90^\circ<x<360^\circ$ nên tập nghiệm của phương  trình trên khoảng $(-90^\circ; 360^\circ)$ là $S=\left\{-75^\circ, 105^\circ, 285^\circ\right\}$.
	 }
\end{ex}


%%%%%%%%%%%%%%%% TỰ LUẬN %%%%%%%%%%%%%%%%

% nguồn: Ok_TuanLuan_Toan11_Lan2
\begin{ex}%[1D1H5-4]%[ma:660b84b6]
	Giải phương  trình $\tan \left(3x-30^\circ \right)=-\dfrac{\sqrt{3}}{3}$.
	\loigiai{ 
		Ta có 
		\begin{eqnarray*}
			&&\tan \left(3x-30^\circ \right)=-\dfrac{\sqrt{3}}{3}\\ &\Leftrightarrow&\tan \left(3x-30^\circ  \right)=\tan \left(-30^\circ \right)\\
			&\Leftrightarrow&x=k60^\circ, k\in  \mathbb{Z}.
		\end{eqnarray*}
		Vậy phương trình có một họ nghiệm  $x=k60^\circ, k\in \mathbb{Z}$.
	 }
\end{ex}

% nguồn: Ok_TuanLuan_Toan11_Lan2
\begin{ex}%[1D1H5-4]%[1D1H5-3]%[ma:2235d695]
	Giải các  phương trình lượng giác sau:
	\begin{multicols}{3}
		\begin{enumerate}
			\item $\cot  4x=\cot\dfrac{2\pi}{7}$;
			 \item $\cot 3x=-2$;
			\item  $\cot\left(2x-10^{\circ}\right)=\dfrac{1}{\sqrt{3}}$.
		\end{enumerate}
	\end{multicols}
	 \loigiai{
		\begin{enumerate}
			\item Ta  có $\cot 4x=\cot\dfrac{2\pi}{7}\Leftrightarrow 4x=\dfrac{2\pi}{7}+k\pi\Leftrightarrow x=\dfrac{\pi}{14}+\dfrac{k\pi}{4}, k\in\mathbb{Z}$.\\
			 Vậy phương trình đã cho có nghiệm là $x=\dfrac{\pi}{14}+\dfrac{k\pi}{4}, k\in\mathbb{Z}$.
			 \item Ta có $\cot 3x=-2=\cot\alpha\Leftrightarrow 3x=\alpha+k\pi\Leftrightarrow x=\dfrac{\alpha}{3}+\dfrac{k\pi}{3}, k\in\mathbb{Z}$.\\
			Vậy phương trình đã cho có nghiệm là $x=\dfrac{\alpha}{3}+\dfrac{k\pi}{3}, k\in\mathbb{Z}$. 
			\item Ta có 
			\begin{eqnarray*}
				&&\cot \left(2x-10^{\circ} \right)=\dfrac{1}{\sqrt{3}}\\ 
				 &\Leftrightarrow& \cot \left(2x-10^{\circ} \right)=\cot 60^{\circ}\\
				&\Leftrightarrow& 2x-10^{\circ}=60^{\circ}+k180^{\circ}\\ &\Leftrightarrow& 2x=70^{\circ}+k180^{\circ}\\ &\Leftrightarrow& x=35^{\circ}+k90^{\circ},  k\in \mathbb{Z}.
			\end{eqnarray*}
			Vậy phương trình đã  cho có nghiệm là $x=35^{\circ}+k90^{\circ}, k\in \mathbb{Z}$.
		\end{enumerate}
	 }
\end{ex}

% nguồn: Ok_TuanLuan_Toan11_Lan2
\begin{ex}%[1D1V5-4]%[ma:e046ad8e]
	Tìm nghiệm của phương trình $2\sin \left(x+40^\circ \right)=\sqrt{3}$ trên khoảng  $\left(-180^\circ;180^\circ \right)$.
	 \loigiai{
		 Ta có
		\begin{eqnarray*}
			2\sin \left(x+40^\circ  \right)=\sqrt{3} &\Leftrightarrow& \sin \left(x+40^\circ \right)=\dfrac{\sqrt{3}}{2}\\
			&\Leftrightarrow& \sin \left(x+40^\circ \right)=sin \left(60^\circ  \right)\\
			&\Leftrightarrow& \hoac{&{x+40^\circ=60^\circ+k360^\circ}  \\
			 &{x+40^\circ=120^\circ+k360^\circ}
			 &} \\
			&\Leftrightarrow& \hoac{&{x=20^\circ+k360^\circ}  \\
			&{x=80^\circ+k360^\circ} 
			 &},\,k\in \mathbb{Z}.
		\end{eqnarray*}
		Theo đề  bài
		\begin{itemize}
			\item  $-180^\circ < 20^\circ+k360^\circ < 180^\circ \Leftrightarrow-\dfrac{5}{9} < k < \dfrac{4}{9} \Rightarrow k=0\Rightarrow x=20^\circ$.
			\item $-180^\circ < 80^\circ+k360^\circ < 180^\circ \Leftrightarrow-\dfrac{13}{18} < k < \dfrac{5}{18}  \Rightarrow k=0 \Rightarrow k=0\Rightarrow x=80^\circ$.
		\end{itemize}
		Vậy  phương trình có hai nghiệm $x=20^\circ$ và $x=80^\circ$.
	 }
\end{ex}
